\begingroup
\sloppy
\section{Unified Mobile Files, Saved Items and Continuation}
\label{sec:pilot-unified-mobile-library}

The unified Pilot mobile application now brings private files, items saved from the television,
research results and subsequent action proposals into one Files surface.  The implementation
reuses Pilot's encrypted Inbox, Private Programme Saves, AION research and Saved Item
Follow-Through authorities.  It does not create another file store or a second saved-item ledger.

\subsection{User Experience}

The Personal space's Files shortcut presents three related collections:
\begin{enumerate}
  \item encrypted attachments that were shared with the active person through the unified Inbox;
  \item products, recipes, destinations, music, learning material and ideas privately claimed
        from shared-television context; and
  \item evidence-backed research, local drafts and separately governed service proposals derived
        from a saved item.
\end{enumerate}

An attachment is fetched only after a fresh signed authorization check, decrypted on the mother
brain and verified against its stored content hash before download.  A television save initially
appears as a short-lived \texttt{awaiting\_private\_claim} item.  The shared screen does not select a
person.  The owner presses ``Save to me'' on a trusted phone, after which the item is isolated to
that person's private library.

For an owned saved item the phone can request research, a shortlist, a task or a learning outline.
Category-appropriate continuations can also prepare shopping, trip or playlist proposals.  Public
evidence links open on the phone so private review can continue without replacing television
playback.

\subsection{Identity and Capability Boundary}

The mobile surface uses the production private identity selected during possession-bound phone
pairing.  The existing household service persona is resolved only inside the mother brain and is
never exposed to the client.  Multi-adult households receive independent saved-item collections;
another adult cannot enumerate, continue or delete an item.  The bridge follows the same safe
legacy-adoption rule as the service-action system: an earlier local-owner record may be adopted
only when exactly one active adult identity exists.

Three narrow signed lease scopes separate the operations:
\begin{itemize}
  \item \texttt{library.read} exposes only that person's saved items and continuation records;
  \item \texttt{library.write} permits private claim and deletion; and
  \item \texttt{library.continue} permits research or preparation from an owned item.
\end{itemize}

Every mobile request is signed by the phone's non-exportable key and checked against its mother-
signed certificate, active identity and short-lived capability lease.  Claiming requires the
current unexpired television-save identifier.  Continuing or deleting requires an already owned
save.  Repeated continuation and deletion requests use content-bound idempotency keys: the exact
retry returns the existing record, while an altered action under the same key fails closed.

\subsection{Evidence and Network Safety}

Research is routed through the provider-independent AION research authority.  Pilot retains a
bounded answer, display label and at most five evidence records.  Only public HTTPS continuation
targets are returned to the phone.  Plain HTTP, localhost, \texttt{.local}, private, loopback,
link-local, multicast and reserved IP targets are removed.  Source photographs and raw audio are
not exposed or retained by this projection.

Research has \texttt{external\_effect=false}.  Opening an evidence page is a user navigation event,
not consent to transact.  A saved item's original confidence describes source identification, not
truth, price accuracy, availability or entitlement.

\subsection{Separate Action Authority}

Shortlist, task and learning continuations create inspectable local drafts only.  Shopping, trip
and playlist continuations create a new Service Execution Hub proposal with its own exact scope,
parameter hash and private approval lifecycle.  Authority is deliberately non-transitive:

\begin{quote}
Saving an item does not authorize research; research does not authorize an account action; and an
account-action proposal does not authorize execution.
\end{quote}

The phone therefore labels a transactional continuation as a separate proposal.  Purchasing,
booking, messaging, calendar changes and provider-account writes remain impossible until their
specific adapter, unchanged approval and execution requirements are independently satisfied.

\subsection{Deletion and Auditability}

The owner can remove a saved item from the active private library.  Pilot records a minimal
idempotent deletion tombstone so a repeated mobile request cannot produce conflicting results.
Historical audit records and any independently approved service proposal remain intact; deleting
the library shortcut cannot erase evidence of an action that was already proposed or executed.

\subsection{Verification and Remaining Qualification}

Automated tests cover private claim, persona isolation, deterministic retries, research
continuation, private-network URL rejection, separate service proposal creation and idempotent
deletion.  Mobile-surface tests cover encrypted file opening, saved-item actions, explicit approval
language and the signed library routes.  At coded closure of \texttt{MOB-0508}, all 464 AION Fabric
tests passed and the optimized Next.js build generated \texttt{/pilot/mobile} successfully.

No purchase, booking, message, calendar change or provider-account write was performed during
verification.  Real-phone qualification must still confirm attachment download, television claim
expiry and reconnection on representative iPhone and Android devices.  Provider execution remains
subject to its own explicit account authorization and verified receipt.
\endgroup
